\section{结论、局限性与未来工作}
\label{sec:conclusion}
\label{sec:limitations}

\paragraph{结论。} 
受 one-shot RLVR~\citep{wang2025oneshot} 的启发，我们提出了一个问题：单次查询是否也足以实现 OPD？答案是肯定的：one-shot OPD 能使学生在数百步内持续提升，并恢复绝大部分全数据 OPD 所带来的收益，跨越不同模型家族和任务领域。数据与算法两个层面给出了统一的答案：OPD 存在数据过载但算法匮乏的问题。一次查询通过其轨迹所到达的状态作用于学生，而单次查询已覆盖全数据 OPD 所访问的大部分状态空间；16 个语义多样化的查询在单领域 OPD 和 MOPD 中均能实现与全数据训练相当的效果。限制运行时长的因素是吸收率：每次更新所吸收的剩余教师-学生差距比例持续下降，无论是一次查询还是全部 17k 次查询，其下降速度一致，因此决定运行时长的是优化器而非训练集。输入只需启动学生的推理过程，这使得数据设计从收集问题转变为选择教师。

\paragraph{局限性。}
状态覆盖是基于全数据回放构建的参考空间进行语义层面的代理度量，它报告的是查询集达到该空间的程度，而非其自身所覆盖的内容，且对每个簇赋予相等权重，无论该簇被访问的频率如何，或仍携带多少教师信号。吸收率的决定因素尚不明确。我们的MOPD运行使用了三个领域，每个领域对应一个教师，因此随着教师数量的增加，多样性结果能扩展到何种程度仍不清楚。

\paragraph{未来工作。}
三个方向依次如下。
\begin{itemize}[topsep=2pt, partopsep=0pt, leftmargin=12pt, itemsep=2pt]
    \item 按状态覆盖选择数据：问题不再是如何收集查询数量，而是这些查询引发哪些状态，因此一旦能够从查询本身估算出状态覆盖，而无需通过完整数据运行来测量，状态覆盖便成为一种选择准则。
    \item 通过在每个token的差距信任区域下，重复使用批量数据进行多轮训练，或根据token所携带的教师信号量对其进行加权，从而提高训练的步骤效率。
    \item 扩展情景，涵盖更多教师和领域，以及更大规模的学生、Agentic 工具使用和长上下文，这些场景下训练数据的收集成本最高。
\end{itemize}